\section{Theory}
\label{sec:theory}

\subsection{The Proportionality Guarantee}

\begin{theorem}[Proportionality]
\label{thm:proportionality}
Let a state have Democratic fraction $d \in (0,1)$, $k$ equal-population
districts, and proportional allocation $k_D = \lfloor dk + 0.5 \rfloor$,
$k_R = k - k_D$.
If the dual-constraint bisection achieves $\mathtt{tpwgts}$ as in
\eqref{eq:tpwgts-formula} exactly, then:
\begin{enumerate}
  \item The R-bloc ($k_R$ districts) has Democratic concentration $= 50\,\%$,
        so Republicans win any competitive R-bloc district.
  \item The D-bloc ($k_D$ districts) has Democratic concentration
        $= \frac{1 - k_R/(2dk)}{k_D/k} \cdot d = \frac{dk - k_R/2}{k_D} > 50\,\%$,
        so Democrats win any competitive D-bloc district.
  \item Recursive bisection within each bloc preserves the majority, so the final
        plan has at least $k_D$ Democratic-winning and $k_R$ Republican-winning districts.
\end{enumerate}
\end{theorem}

\begin{proof}
Item~(1): Democratic concentration in R-bloc $= \frac{(1-\tau_D) D}{(k_R/k) P}
= \frac{k_R/(2dk) \cdot dP}{(k_R/k) P} = \frac{1}{2}$. \checkmark

Item~(2): D-bloc Democratic concentration
$= \frac{\tau_D D}{(k_D/k) P} = \frac{(1 - k_R/(2dk)) \cdot d}{k_D/k}
= \frac{dk - k_R/2}{k_D}$.
Since $k_D = dk$ rounded up (under HH), $k_D \geq dk$ and $k_R \leq (1-d)k$.
Therefore: $\frac{dk - k_R/2}{k_D} \geq \frac{dk - (1-d)k/2}{k_D}
= \frac{k(3d-1)/2}{k_D}$.
For $d > 1/3$ (D has at least one seat), this exceeds $1/2$ when $k_D \leq k(3d-1)/1$.
In the two-party case with $d > 0.5$ (D plurality), the concentration is always $> 50\,\%$. \checkmark

Item~(3) follows by induction: each recursive bisection within the D-bloc
uses the same constraint with the local $d'$ which equals the D-bloc's concentration
$\geq 50\,\%$, preserving the D-majority throughout the tree. \qed
\end{proof}

\begin{remark}
For $d < 0.5$ (R-plurality state), $k_D < k/2$. The formula still applies,
but the R-bloc gets the majority of districts. The guarantee is symmetric:
each party wins its proportional allocation.
\end{remark}

\subsection{The Partisan Lorenz Feasibility Condition}

Not all states can achieve the $\mathtt{tpwgts}$ targets under the geographic
contiguity constraint. The partisan Lorenz curve determines feasibility.

\begin{definition}[Partisan Lorenz curve]
\label{def:lorenz}
Sort census tracts by Democratic fraction $d_v / p_v$, from least to most
Democratic (R-heavy first). Define $L_D : [0,1] \to [0,1]$ by:
\[
  L_D(x) = \frac{\sum_{v : \mathrm{rank}_R(v) \leq x|V|} d_v}{D}
\]
where tracts are ordered from least to most Democratic.
$L_D(x)$ is the fraction of all Democratic votes in the $x$-fraction least-Democratic tracts.
\end{definition}

\begin{theorem}[Lorenz feasibility]
\label{thm:lorenz-feasibility}
The proportional target $\tau_D^* = 1 - k_R/(2dk)$ is geometrically achievable
(without contiguity constraint) if and only if:
\[
  L_D(k_R/k) \leq \frac{k_R}{2dk}
\]
i.e., the least-Democratic $k_R/k$ fraction of the population contains at most
the required minimum of Democratic votes for the R-bloc.
\end{theorem}

\begin{proof}
$L_D(k_R/k)$ is the maximum Democratic fraction achievable in the R-bloc under
non-contiguous selection of $k_R/k$ of the population (taking the most
Republican-friendly tracts). If this maximum is $\leq k_R/(2dk)$, the target
is achievable; if it exceeds $k_R/(2dk)$, even the most Republican-friendly
geography contains too many Democrats for a proportional R-bloc. \qed
\end{proof}

\begin{definition}[Natural proportional ratio]
The \emph{natural proportional ratio} $p^*$ is the population fraction at which
the partisan Lorenz curve crosses the proportionality target:
\[
  p^* = \inf\{x : L_D(x) \geq k_R/(2dk)\}
\]
When $p^* \leq k_R/k$, proportionality is feasible. When $p^* > k_R/k$,
more population is needed than the R-bloc's share to find enough R-friendly tracts.
\end{definition}

\subsection{The Minimum Partisan Signal}

\begin{definition}[Partisan signal strength]
The \emph{partisan signal strength} $\sigma$ of a redistricting algorithm is the
$L^1$ distance between its $\mathtt{tpwgts}[1]$ (D-votes target for left partition)
and the neutral target $k_D/k$ (equal Democratic distribution):
\[
  \sigma = \left|\tau_D - \frac{k_D}{k}\right| = \left|1 - \frac{k_R}{2dk} - \frac{k_D}{k}\right|
\]
\end{definition}

$\sigma = 0$ means no partisan signal is required: the algorithm is neutral
and achieves proportionality by geography alone.
$\sigma > 0$ measures how far the algorithm must ``pull'' Democratic votes
from the neutral geographic distribution to achieve proportionality.

\begin{proposition}
For a state with uniform partisan geography (each tract has Democratic fraction
exactly $d$), $\sigma = 0$ and proportionality is achieved by any balanced bisection.
For a state where all Democratic voters are concentrated in one geographic cluster,
$\sigma$ is maximised and geographic contiguity may prevent proportionality entirely.
\end{proposition}
